\chapter{CODE}
\label{ch:code}

\section{ESP32 Main Controller Code}
\hspace{2em}The following C++ code snippet represents the core firmware execution loop implemented in the production VibeGuard firmware (\texttt{VibeGuard\_Autotune\_Synchronized.ino}). It implements real-time 800\,Hz sensor acquisition over SPI, in-place zero-mean DC cancellation, 3-axis Vector RMS extraction, adaptive threshold evaluation, and $K$-of-$M$ persistence filtering.

\begin{lstlisting}[language=C++, caption={VibeGuard Core Real-Time Acquisition and DSP Firmware Loop}, label={lst:vibeguard_core}]
void loop() {
    // 1. Acquire 256 samples @ 800 Hz (320 ms window)
    AccelGData g_sample;
    for (size_t i = 0; i < 256; ++i) {
        accel.readG(g_sample);
        x_buf[i] = g_sample.x;
        y_buf[i] = g_sample.y;
        z_buf[i] = g_sample.z;
        delayMicroseconds(1250); // Deterministic 800 Hz pacing
    }

    // 2. Remove DC gravity offset and static mounting tilt
    remove_dc(x_buf, 256);
    remove_dc(y_buf, 256);
    remove_dc(z_buf, 256);

    // 3. Compute Triaxial Vector RMS magnitude
    float v_rms = compute_vector_rms(x_buf, y_buf, z_buf, 256);

    // 4. Update Warning and Alarm persistence queues (K=4, M=5)
    bool is_warn = (v_rms >= g_profile.warning_threshold_g);
    bool is_alarm = (v_rms >= g_profile.alarm_threshold_g);
    
    update_persistence(is_warn, is_alarm);
    evaluate_condition_state();

    // 5. Update physical RGB LED and transmit telemetry
    applyIndicator();
    transmit_json_telemetry(v_rms);
}
\end{lstlisting}

\section{Multi-Stage Physical RGB LED Indication Logic}
\hspace{2em}Listing~\ref{lst:led_actuation} details the multi-color RGB LED actuation routine responsible for visually rendering the physical operating states (Normal, Warning, Alarm, Calibration, Fault).

\begin{lstlisting}[language=C++, caption={Multi-Stage Physical RGB LED Indicator Actuation}, label={lst:led_actuation}]
void applyIndicator() {
    // 1. Hardware or Sensor Acquisition Fault -> Purple / Magenta
    if (!g_sensor_ready || g_committed_state == ConditionState::SENSOR_FAULT) {
        setRGB(true, false, true); // Red + Blue = Purple
        return;
    }

    // 2. Active Baseline Calibration (8.0 s Median + MAD) -> Solid Blue
    if (g_vibeEngine.getCalibrationState() != CAL_IDLE) {
        setRGB(false, false, true); // Solid Blue
        return;
    }

    // 3. Physical Operating States
    switch (g_committed_state) {
        case ConditionState::ALARM: {
            // Stage 2/3 Severe Unbalance -> Fast Red Strobe @ 6.25 Hz
            bool strobe = (millis() % 160) < 80;
            setRGB(strobe, false, false);
            break;
        }
        case ConditionState::WARNING: 
            // Stage 1 Moderate Unbalance -> Solid Yellow (Red + Green)
            setRGB(true, true, false); 
            break;
        case ConditionState::NORMAL:  
            // Stage 0 Healthy Baseline -> Solid Green
            setRGB(false, true, false); 
            break;
        default: 
            setRGB(false, false, true); 
            break;
    }
}
\end{lstlisting}

\section{Substrate-Invariant Rotordynamic DSP Engine Implementation}
\hspace{2em}Listing~\ref{lst:engine_dsp} illustrates the deep research rotordynamic DSP pipeline implemented in \texttt{VibeGuardEngine.h}, showcasing the biquad low-pass cascade, autocorrelation shaft speed tracker, and the dimensionless condition index $J(t)$.

\begin{lstlisting}[language=C++, caption={Substrate-Invariant Rotordynamic Metric $J(t)$ Extraction}, label={lst:engine_dsp}]
// Biquad Direct Form II 4th-Order Butterworth (25 Hz cutoff @ 800 Hz ODR)
float processBiquad(float in, BiquadState& s, const BiquadCoeffs& c) {
    float w = in - c.a1 * s.w1 - c.a2 * s.w2;
    float out = c.b0 * w + c.b1 * s.w1 + c.b2 * s.w2;
    s.w2 = s.w1;
    s.w1 = w;
    return out;
}

// Bounded Autocorrelation Shaft Speed Tracker (7.5 - 14.0 Hz)
float estimateShaftSpeedHz(const float* x, int N) {
    float best_r = -1.0f;
    int best_lag = 80; // 10.0 Hz nominal
    for (int k = 55; k <= 105; ++k) {
        float r = 0.0f;
        for (int i = 0; i < N - k; ++i) r += x[i] * x[i + k];
        if (r > best_r) { best_r = r; best_lag = k; }
    }
    return 800.0f / static_cast<float>(best_lag);
}

// Dimensionless Condition Metric J(t) with Speed-Sag Compensation
float computeMetricJ(float Exy_LF, float Rx, float Ry, float CF, float f1) {
    float c_sag = sq(cal_.f1_0 / f1);
    float term_energy = (Exy_LF / cal_.Exy_LF_0) - 1.0f;
    float term_product = ((Rx / cal_.Rx_0) * (Ry / cal_.Ry_0) * (CF / cal_.CF_0)) - 1.0f;
    return c_sag * (0.55f * term_energy + 0.45f * term_product);
}
\end{lstlisting}

\section{Code Walkthrough and Implementation Details}
\hspace{2em}The implementation highlighted in Listings~\ref{lst:vibeguard_core}, \ref{lst:led_actuation}, and \ref{lst:engine_dsp} incorporates several embedded software principles:

\begin{enumerate}
    \item \textbf{Deterministic Hardware Pacing:} Sampling is strictly regulated at $1250\,\mu\text{s}$ intervals using the high-resolution hardware microsecond timer (\texttt{delayMicroseconds}), guaranteeing an exact 800\,Hz Output Data Rate without timing jitter.
    \item \textbf{In-Place Signal Preprocessing:} The static DC components (Earth's $1.0\,g$ gravity bias and physical sensor tilt) are computed and subtracted directly within the input sample arrays using \texttt{remove\_dc()}, avoiding duplicate memory allocations.
    \item \textbf{Direct Form II Biquad Filtering:} Cascaded 2nd-order stages suppress gear tooth mesh impacts ($150\text{--}350\,\text{Hz}$) by over $36\,\text{dB}$ while maintaining minimal state register overhead (two delay elements per biquad).
    \item \textbf{Substrate-Invariant Ratio Metric:} Computing $J(t)$ as normalized quotients guarantees mathematical homogeneity of degree zero, ensuring identical threshold sensitivity regardless of whether the rig rests on wood, acrylic, or steel foundations.
    \item \textbf{Dual-Stage Persistence Filtering:} Independent circular queues track both warning and alarm violations over sliding 5-window depths ($M=5$), preventing transient shocks from provoking premature alarm states.
    \item \textbf{Non-Blocking Strobe Generation:} The 6.25\,Hz Fast Red Strobe utilizes modulo arithmetic on \texttt{millis()} rather than blocking \texttt{delay()} calls, preserving real-time telemetry streaming and sensor sampling continuity.
    \item \textbf{Zero-Heap Allocation in Steady State:} All sample buffers and filter states are statically allocated in internal SRAM, eliminating dynamic memory fragmentation and guaranteeing long-term continuous uptime.
\end{enumerate}
\clearpage
